% !TeX root = ../../main.tex
\subsection{Inverse domains and characteristic geometry}
\label{subsec:b:domains}

We use the entire Bessel series and the critical data of
\eqref{eq:s06:critical-data}. Since $w''(x_*)<0$, the implicit
function theorem defines, for real $\theta$ sufficiently close to zero,
\begin{equation}\label{eq:b:critical-data}
 w'(x_{*,\theta})+\theta=0,\qquad
 x_{*,0}=x_*,\qquad R_\theta=w_\theta(x_{*,\theta}).
\end{equation}
These functions are real analytic. We will always shrink their
parameter interval when necessary. For complex disks write
$\mathbb D(c,r)=\{u\in\mathbb C:|u-c|<r\}$, where $r>0$.

The construction has two separate steps. A derivative estimate makes
the inverse single valued on a convex domain. A difference-quotient
estimate then proves the full radius needed by the kernels; the latter
does not follow from univalence alone.

Only in this subsection, for integers $k\geq0$, put
\[
 M_k(u)=\sum_{a\geq0}\frac{(-u)^a}{a!(a+k)!}.
\]
Thus $M_0=w'$ and $w^{(k+1)}=(-1)^kM_k$.

\begin{lmm}[Angular bounds]\label{lmm:b:angular}
One has $x_*<1.446$. For $0\leq r\leq x_*$, $|u|=r$ and
every integer $k\geq0$,
\begin{equation}\label{eq:b:angular}
 \operatorname{Re}M_k(u)\geq M_k(r).
\end{equation}
If $r>0$ and $u\ne r$, the inequality is strict. For $k\geq1$
and $0\leq r\leq1.447$ one also has
\begin{equation}\label{eq:b:derivative-signs}
 0<\frac1{k!}-\frac r{(k+1)!}
                  \leq M_k(r)\leq\frac1{k!}.
\end{equation}
\end{lmm}
\begin{prf}
The even partial sum through degree six of $M_0(1.446)$ is less
than $-8\cdot10^{-5}$. Its tail is alternating with decreasing
terms, so $M_0$ has a positive zero before $1.446$. This bounds
the first zero, which is $x_*$.

For $k=0$, write $u=re^{i\varphi}$ and $c=\cos\varphi$.
Define the Chebyshev polynomials by
$T_a(\cos\varphi)=\cos(a\varphi)$, and put
$\mathcal U_r(c)=\sum_{a\geq0}(-r)^aT_a(c)/(a!)^2$.
The finite sine-quotient formula gives
$|T_a'|\leq a^2$ and
$|T_a''|\leq a^2(a^2-1)/3$ on $[-1,1]$.
Consequently, for $0<r\leq1.45$,
\[
 r^{-2}\mathcal U_r''(c)
 =1-\frac{2r c}{3}
       +r^2\left(\frac{c^2}{6}-\frac1{36}\right)+E_r(c),
 \qquad |E_r(c)|<0.049.
\]
For completeness, the terms of the absolute tail at $a=5,6$ are
less than $0.042343$ and $0.003582$. After $a=6$ their consecutive
ratio is at most
$1.45(a+2)/(a(a^2-1))<0.056$; this proves the stated tail bound.
The displayed polynomial decreases in $c$ on $[-1,1]$, since
its $c$ derivative is $r(rc-2)/3<0$. At $c=1$ it also decreases for
$0<r\leq1.45$, and its minimum is $937/2880>0.325$.
Thus $\mathcal U_r''>0$. Moreover
$\mathcal U_r'(1)=-rM_0(r)\leq0$ for $r\leq x_*$.
It follows that $\mathcal U_r(c)>\mathcal U_r(1)$ when $c<1$.

For $k\geq1$, the first odd term of
$\operatorname{Re}M_k(re^{i\varphi})-M_k(r)$ is
$r(1-\cos\varphi)/(k+1)!$; every other odd term is nonnegative.
Using $1-\cos(a\varphi)\leq a^2(1-\cos\varphi)$, the sum of the
possible negative even terms divided by that first term is at most
\[
 \frac{2r}{3}+\frac{r^3}{90}
       +\sum_{\substack{a\geq6\\a\ \text{even}}}
                         \frac{2a^2r^{a-1}}{a!(a+1)!}<1
                       \qquad(0\leq r\leq1.447).
\]
Indeed the last sum is less than $0.0002$: its first term is
$r^5/50400$ and its successive ratios are at most
$1.447^2/(a^2(a+1)(a+3))<0.001$.
This also proves strictness. Finally, the alternating terms of
$M_k(r)$ decrease from the first term for the range in question;
the first two partial sums give \eqref{eq:b:derivative-signs}.
\end{prf}

\begin{lmm}[A common convex inverse domain]\label{lmm:b:convex-domain}
There are $a_*\in(0,x_*/2)$ and $\theta_*>0$ such that, for
$0\leq\theta\leq\theta_*$, the disk
\begin{equation}\label{eq:b:omega}
 \Omega_\theta=\mathbb D(a_*,x_{*,\theta}-a_*)
\end{equation}
contains zero and satisfies
$\operatorname{Re}w_\theta'(u)>0$ on its closure except at
$u=x_{*,\theta}$. For every $\sigma\geq0$, $w_{\theta+\sigma}$
is injective on $\overline\Omega_\theta$ and its inverse has positive
derivative real part in its image.
\end{lmm}
\begin{prf}
Choose a positive $a_*$ as small as necessary below.
At $\theta=0$, the closed disk lies in $|u|\leq x_*$, and
Lemma~\ref{lmm:b:angular} gives strict positivity except at $x_*$.
Near the moving right endpoint write
$u=x_{*,\theta}-A+iB$, with $A\geq0$. The disk condition gives
$B^2\leq2(x_{*,\theta}-a_*)A-A^2$. The Bessel equation, obtained
also by differentiating the series, is
\begin{equation}\label{eq:b:bessel-ode}
 u w'''(u)+w''(u)+w'(u)=0.
\end{equation}
At the endpoint it gives
\[
 -w''(x_{*,\theta})
 -(x_{*,\theta}-a_*)w'''(x_{*,\theta})
                      =a_*w'''(x_{*,\theta})-\theta.
\]
The right side is positive when $\theta$ is small, since
$w'''(x_*)=-w''(x_*)/x_*>0$. Taylor expansion therefore bounds
$\operatorname{Re}w_\theta'(u)$ below by a positive multiple of
$A$, on a fixed small endpoint neighborhood. The remainder is
$O(A|u-x_{*,\theta}|)$, because $|u-x_{*,\theta}|^2\leq C A$.
On the rest of the closed disk, strict positivity at $\theta=0$
and compactness give the conclusion after reducing $\theta_*$.

Adding $\sigma$ preserves the positive real part. For distinct
$u,v$ in the closed disk,
\[
 \frac{w_{\theta+\sigma}(u)-w_{\theta+\sigma}(v)}{u-v}
 =\int_0^1w_{\theta+\sigma}'(v+t(u-v))\,\dd t
\]
has positive real part, since the open segment lies in the disk.
This proves injectivity even on its closure. The reciprocal of a
number with positive real part also has positive real part, proving
the assertion about the inverse.
\end{prf}

\begin{prp}[The full inverse-radius budget]\label{prp:b:inverse-budget}
The constants in Lemma~\ref{lmm:b:convex-domain} can be chosen so that,
with $D_\theta=w_\theta(\Omega_\theta)$, every real
$0\leq x<R_\theta$ satisfies
\begin{equation}\label{eq:b:inverse-budget}
 \operatorname{dist}(x,\partial D_\theta)=R_\theta-x.
\end{equation}
The closest boundary point is uniquely $R_\theta$. The same budget
and unique contact hold on a sufficiently small fixed negative real
interval, uniformly in $\theta$. The inverse
$v_\theta$ extends across every other point of $|x|=R_\theta$.
At the remaining point it has the convergent local expansion
\begin{equation}\label{eq:b:morse-inverse}
 v_\theta(R_\theta-S)
 =x_{*,\theta}
       -\sqrt{\frac{2S}{-w''(x_{*,\theta})}}+O(S),
\end{equation}
on sectors $|\arg S|<\pi-\gamma$, for each fixed $\gamma>0$.
The coefficients and local bounds are uniform for small
$0\leq\theta\leq\theta_*$. Square roots in this appendix use the
branch positive on the positive real axis.
\end{prp}
\begin{prf}
First let $x=w_\theta(v)$ with $a_*\leq v<x_{*,\theta}$ and put
$r=x_{*,\theta}-v$. The disk $\mathbb D(v,r)$ is contained in
$\Omega_\theta$ and tangent to its boundary at the critical point.
The difference quotient on its boundary is
\[
 \frac{w_\theta(v+s)-w_\theta(v)}s
 =w_\theta'(v)+\sum_{k\geq1}
                       \frac{(-1)^kM_k(v)s^k}{(k+1)!}.
\]
Its real part is minimized at $s=r$. To check this directly, divide
the possible negative even angular differences by
$M_1(v)r(1-\cos\varphi)/2$.
By \eqref{eq:b:derivative-signs} their sum is at most
\[
 \frac{2x_{*,\theta}}3+\frac{x_{*,\theta}^3}{90}+0.0002<1.
\]
Here $r/(1-v/2)\leq x_{*,\theta}$, and we shrink $\theta_*$ so
that $x_{*,\theta}<1.447$. This is the same even-term estimate as
in Lemma~\ref{lmm:b:angular}, with the first odd term retained.
It is strict unless $s=r$. Thus
\[
 |w_\theta(v+s)-x|\geq R_\theta-x\quad(|s|=r),
\]
with equality only at the tangency point. Univalence then shows that
the image contains $\mathbb D(x,R_\theta-x)$, with no other
boundary contact.

We also need this conclusion at $v=0$. For $a_*=\theta=0$,
\eqref{eq:s06:circle-bound} gives
$|w(u)|>R$ on $|u|=x_*$ except at $u=x_*$. Away from that point
the strict margin persists under a sufficiently small positive
choice of $a_*$ and then a small $\theta$. Near the moving endpoint,
the circle has $A=B^2/(2(x_{*,\theta}-a_*))+O(B^4)$, so
\[
 \operatorname{Re}(w_\theta(u)-R_\theta)
             =-\tfrac12w''(x_{*,\theta})B^2+O(|B|^3)>0.
\]
This supplies the missing local strict margin at $x=0$.
For $0\leq v\leq a_*$, fix a boundary image $z$ and observe that
$|z-x|^2-(R_\theta-x)^2$ is affine in the real variable $x$.
It is nonnegative at $x=0$ and at $x=w_\theta(a_*)$, and strictly
positive there unless $z=R_\theta$. This completes
\eqref{eq:b:inverse-budget}.

The negative extension follows from the same affine expression.
Near the critical boundary image we have just proved
$\operatorname{Re}z>R_\theta$, so negative $x$ only improves it.
Away from that point the strictly positive compact margin absorbs a
sufficiently small negative $x$, uniformly in $\theta$. Such centers
remain in $D_\theta$ because zero is an interior point.

The closed-disk injectivity proved above identifies the boundary
image without self-intersections. At its regular points the inverse
function theorem gives continuation. At the critical point,
$R_\theta-w_\theta(u)$ has a double zero with coefficient
$-w''(x_{*,\theta})/2>0$. Taking the analytic square root of the
nonvanishing factor and inverting the resulting local coordinate
proves \eqref{eq:b:morse-inverse}, with its uniform sector bounds.
\end{prf}

For the characteristic geometry put $S=R_\theta-x$. We use the
following two-distance cones. For $0<\gamma<\pi/2$, let
$\mathcal C_\gamma$ consist of pairs $(S,y)$ with $S\ne0$,
$S\pm y\notin(-\infty,0]$, for which
\begin{equation}\label{eq:b:cone}
 \max\{0,\arg(S-y),\arg(S+y)\}
 -\min\{0,\arg(S-y),\arg(S+y)\}<\pi-\gamma.
\end{equation}
Arguments are principal. Positive linear combinations of the two
distances and of positive real numbers stay in this cone. Their
moduli are bounded below by a constant depending on $\gamma$ times
the sum of the moduli. In particular $|y|\leq C_\gamma|S|$.

\begin{lmm}[Shifts and characteristics]\label{lmm:b:cone-transport}
Fix $0<\gamma<\pi/4$. Kernel shifts $(x,y)\mapsto(x+c y,t y)$, where $c$ is real,
$0\leq t\leq1$ and $|c|+t\leq1$, preserve
$\mathcal C_\gamma$. For $(S,y)\in\mathcal C_{2\gamma}$ with
$|S|$ sufficiently small, $0\leq\theta\leq\theta_*/2$ and
$0\leq\sigma\leq\sigma_0$, where $0<\sigma_0\leq\theta_*/2$ is small and fixed,
put $v=v_\theta(x)$ and
$S_\sigma=R_{\theta+\sigma}-(x+\sigma v)$. Then
$(S_\sigma,y)\in\mathcal C_\gamma$ and
\begin{equation}\label{eq:b:transport-distances}
 \begin{aligned}
 |S_\sigma|&\asymp |S|+\sigma\sqrt{|S|}+\sigma^2,\\
 |S_\sigma\pm y|&\geq c_\gamma
       \bigl(|S\pm y|+\sigma\sqrt{|S|}+\sigma^2\bigr).
 \end{aligned}
\end{equation}
These estimates are uniform in the stated ranges. On the full
convex image, characteristic transport also never decreases the
distance from its center to the boundary.
\end{lmm}
\begin{prf}
After a kernel shift, each new distance is a convex combination
of $S-y$ and $S+y$, since $|c|+t\leq1$. This proves the first
assertion and its modulus comparison.
Differentiation at the critical point gives
\[
 R_\theta'=x_{*,\theta},\qquad
 x_{*,\theta}'=-1/w''(x_{*,\theta})>0.
\]
Consequently the exact distance identity is
\begin{equation}\label{eq:b:transport-distance-identity}
 S_\sigma=S+\sigma(x_{*,\theta}-v)
       +\int_0^\sigma(x_{*,\theta+t}-x_{*,\theta})\,\dd t.
\end{equation}
The last term is positive real and comparable to $\sigma^2$.
By \eqref{eq:b:morse-inverse}, the middle coefficient equals a
positive multiple of $\sqrt S$, with relative error $O(\sqrt{|S|})$.
The argument of $\sqrt S$ lies between zero and that of $S$.
The angular margin between $\mathcal C_{2\gamma}$ and
$\mathcal C_\gamma$ absorbs this error, uniformly after reducing
the small constants. Positive-cone addition in
\eqref{eq:b:transport-distance-identity} proves
\eqref{eq:b:transport-distances}.

For the last assertion, use any $u,v\in\overline\Omega_\theta$,
$u\ne v$, and set
$A=\int_0^1w_\theta'(v+t(u-v))\,\dd t$.
Its real part is positive, and
\[
 |w_{\theta+\sigma}(u)-w_{\theta+\sigma}(v)|
              =|u-v|\,|A+\sigma|\geq|u-v|\,|A|.
\]
Taking the infimum over boundary $u$ proves the claim. It is strict
for $\sigma>0$. This argument applies to every interior complex
center, not only the real centers in \eqref{eq:b:inverse-budget}.
\end{prf}

\begin{lmm}[Disks on the negative axis]\label{lmm:b:negative-disks}
For $u>0$, $w$ is univalent on $\mathbb D(-u,x_*)$, and
\begin{equation}\label{eq:b:negative-radius}
 \operatorname{dist}\bigl(w(-u),
                   \partial w(\mathbb D(-u,x_*))\bigr)
 \geq w(x_*-u)-w(-u)>R.
\end{equation}
The strict margins are uniform when $u$ ranges over a compact
subinterval of $(0,\infty)$.
\end{lmm}
\begin{prf}
The entire Taylor expansion in $u$ and
Lemma~\ref{lmm:b:angular} give, for $|s|\leq x_*$,
\[
 \operatorname{Re}w'(-u+s)
 =\sum_{k\geq0}\frac{u^k}{k!}\operatorname{Re}M_k(s)
 \geq\sum_{k\geq0}\frac{u^k}{k!}M_k(|s|)
 =w'(-u+|s|)>0.
\]
The segment argument proves univalence. On $|s|=x_*$, integrate
this inequality along the radial segment to obtain the first bound
in \eqref{eq:b:negative-radius}. Finally $w'$ is strictly decreasing
on $(-\infty,x_*]$: for negative arguments this follows from its
positive coefficient series, and for positive arguments it follows
from \eqref{eq:b:derivative-signs}. Therefore
\[
 w(x_*-u)-w(-u)=\int_0^{x_*}w'(t-u)\,\dd t
                           >\int_0^{x_*}w'(t)\,\dd t=R.
\]
Continuity on compact intervals proves uniformity.
\end{prf}
